\chapter[Why a Urinal Can Be Museum Art]{Why an Ordinary Urinal Can Enter an Art Museum}
\section{A Ceramic Object Facing Away}
First, pick up an object. This time, promise not to frown.

It is a white ceramic urinal, shaped like a flattened shell, flat behind and gently curved in front. Building-supply stores carry rows of similar objects whose purpose is so obvious that even writing "urinal" in a book title gives us pause.

This one differs. Instead of mounted upright, it lies flat on a pedestal like an overturned helmet. Black paint on its downward-facing side reads "R. Mutt 1917." No plumbing is attached or ever will be. Gallery spotlights shine above it; a label stands nearby. Visitors queue beyond glass, frowning, suppressing smiles, taking photographs, or lingering a long time.

This is Marcel Duchamp's Fountain, visible in major museums in London, Paris, and San Francisco, or rather in replicas. The 1917 original disappeared; Duchamp authorized an edition in 1964, now treated as top-tier collections. A few-dollar hardware item, reoriented and signed, became one of twentieth-century art history's most discussed works.

Notice your response: "That's art?", "Cool," or "A scam." The response itself is crucial to the work. What did the urinal do to enter the museum, and how did its admission change our understanding of art?

\section{1917: An Exhibition That Could Reject Nobody}
Enter a scene.

April 1917, New York. European trench reports fill newspapers. The city gathers exiled artists, merchants, and refugees amid strange excitement: the Old World burns; the New thinks anything possible.

Artists establish the Society of Independent Artists, rejecting academic juries that decide which paintings deserve walls and which deserve no viewing. Their generous rule: no jury, no prizes; pay six dollars and every submitted work is accepted and displayed.

Remember this promise of fairness. The chapter's drama rests on it.

On April 10, the exhibition opens in a great hall. Over two thousand works by more than a thousand artists fill it. Equality dictates alphabetical hanging, regardless of fame, placing a master's oil beside an unknown beginner's exercise. Crowds wander, critics take notes: an evening sincerely believing art belongs to everyone. Yet nobody finds the white urinal signed R. Mutt.

One director is the thirty-year-old Frenchman Duchamp, already known for Nude Descending a Staircase, whose overlapping color shapes had drawn American outrage and praise. Early in April, organizers received a white urinal lying flat, signed R. Mutt, sent from Philadelphia, titled Fountain. The full six-dollar fee accompanied it.

Rules required acceptance. Organizers hesitated: no alteration beyond an unfamiliar signature. Was an artwork actually submitted? Later accounts describe heated pre-opening argument. Ultimately it was removed and not exhibited. An exhibition defined by no jury held a judgment on its opening eve.

Duchamp calmly resigned from the board, without argument or statement. The gesture said enough: you broke your own rule. Days later the rejected object reached Alfred Stieglitz's studio for its famous photograph, white ceramic lit before a hanging painting, the urinal lying on the floor. Never entering the exhibition, it entered history through a photograph.

Who was R. Mutt? Almost nobody knew. Later consensus treated it as Duchamp's pseudonym, perhaps parodying plumbing supplier Mott, perhaps drawn from the comic Mutt and Jeff. Duchamp never fixed one answer. An untraceable signature on an object not made by its signer forecast the entire affair.

\section{Do Not Choose a Side Yet}
State the conflict before answering.

Organizers' reasoning sounds reasonable: surely artists must make art? Buy plumbing, sign falsely, demand exhibition, and every plumber becomes a sculptor, every warehouse a museum. Why need artists? What would paying visitors see but a public joke at their expense?

Duchamp's reasoning is equally firm. The written rules promise no jury and display for payment; privately judging an disliked object reveals unpreparedness for unfamiliar art, not opposition to bad art. Who required handmade creation? Choosing, naming, and changing circumstances can create too.

The dispute reaches three deeper questions:

First, \textbf{what determines a work's identity}: ceramic or bronze, three months' craft or three minutes' shopping, beauty, location, or the person placing it there?

Second, \textbf{who may declare something art}: artist, jury, museum, paying public, majority applause? The 1917 exhibition answered everyone until a urinal made it retract.

Third, \textbf{does intention count}? Is a casual prank art, or a deliberate challenge more so? Depending on someone's private thought seems stranger than the urinal itself.

Choose provisionally: on that committee, facing rule and urinal, would you display it?

\section{A Room of Voices and Silent Stones}
Give the room's positions full strength, then examine comparable problems elsewhere.

\textbf{First: speak fully for the organizers.}

Calling rejection simply conservative stupidity is unfair. Three defenses deserve attention. Craft promises years of training in drawing, anatomy, perspective, color; treating shopping as creation publicly dismisses that labor. Museums filter rather than store: walls and audiences' time are finite, and abandoning selection rewards opportunistic spectacle over seriousness. Finally, unrestricted self-declaration can legalize scams: critics become ignorant by definition while shameless claimants remain invulnerable. These arguments return whenever controversial art reaches the news.

\textbf{Second: Duchamp and his friends.}

Art's value includes ideas, not just hands. Architects need not lay bricks, composers need not make instruments; why must visual artists touch all materials themselves? Selecting, reorienting, naming, and submitting are a creative sequence.

\textbf{Third: those who made it.}

The urinal came from a ceramic factory, through mixing, molding, glazing, firing by skilled workers. They really made it, yet art history forgets their names. Glory goes to selection, silence to manufacture. While debating Duchamp's authorship, people rarely invite these qualified makers into the room.

\textbf{Fourth: later collectors and museums.}

Victory produces the most dramatic reversal. Rejected then, the object later enters major collections and every textbook. Institutions enthrone their challenger; trouble becomes canon. Does rebellion retain meaning when it becomes orthodoxy?

Elsewhere we find that selecting and placing objects into art long preceded Duchamp.

Chinese literati admired stones at least from the Song dynasty: choosing natural forms from Lake Tai or mountains, leaving them uncarved, adding stands, names, poems. Selection, placement, naming supply art where humans did not make the stone, remarkably resembling Duchamp. Mi Fu supposedly arranged his clothes and bowed to an unusual rock as "Elder Brother Stone." Whether factual is uncertain; the tale's survival shows an unworked stone as spiritual object was entirely intelligible.

Japanese tea culture often values uneven, roughly glazed bowls over exquisite official wares. Famous Ido bowls began as Korean peasants' everyday rice bowls, then tea practitioners selected, named, and transmitted them at values beyond elaborate tribute goods. They chose ordinary utensils and transformed context through ritual into aesthetic objects.

An opposite movement is more revealing. In many West African societies, masks function in dance, drumming, seasons, and communal ritual, not as wall objects. Removed into European museum cases, they changed from things doing work to sculptures being watched. Context transformed meaning. Many arrived under colonial conditions, making their museum histories matters requiring narration and reckoning. Context involves power as well as aesthetics.

Duchamp's novelty was not selection itself but selection as a weapon: others selected to house beauty; he selected to question institutions.

\section[Three Questions about Art]{A Bridge for Thought: Three Questions within "Is It Art?"}
"That is not art" and "great art" often talk past each other because three questions masquerade as one. Separate them with this portable tool.

\textbf{First, classification: does it belong to art?}

Should it enter the category for artistic discussion, exhibition, history? Admission can be broad enough for urinals, stones, bowls, and masks.

\textbf{Second, evaluation: is it good art?}

Is it profound, new, worth time? Evaluation needs standards and disagreement. Anything may become art without every artwork being equally good. A bowl can qualify yet be mediocre.

\textbf{Third, interpretation: why is it here, and what is intended?}

Whom does it mock, commemorate, challenge? Different intentions and contexts transform the same thing's meaning.

Try Fountain. Defenders primarily need to win classification; critics want evaluation but fire classification's ammunition, declaring it unworthy of the name. They are fighting different battles.

Try everyday examples. Anonymous desk carving can be studied as folk graffiti while judged crude and harmful. Damage to public property raises a fourth independent question, legality; do not use art as its shield. A grandmother's repeatedly repaired blue-rimmed bowl qualifies in an ethnographic museum; its craft may be plain, while its meaning carries decades of family meals, more important than glaze.

Before siding with "does this deserve art's name?", ask which question is disputed. Most confusion dissolves, leaving the genuinely worthwhile disagreement.

\section{A Brief Defense of Mr. Mutt}
Read primary evidence. After Fountain's removal, a small magazine produced by Duchamp and friends responded.

The Blind Man lasted only two issues. Its second, May 1917, largely defended the rejected work, centered on the unsigned "The Richard Mutt Case," generally attributed to Duchamp's circle. Its famous sentence means:

\begin{quote}
Whether Mr. Mutt made the urinal himself does not matter. He chose it.
\end{quote}
\begin{quote}
(From The Blind Man, no. 2, May 1917, New York; a public-domain publication originally in English.)
\end{quote}
"He chose it" bears the defense's weight. It continues broadly: selection from daily life, new placement, title, and viewpoint remove practical use and create a new thought for the object.

It also half-jokingly answers "mere plumbing": America's contributions to world art are plumbing and bridges. The joke has substance: admired modern life rests not solely on artists' hands but on disparaged industrial products.

The strategy never claims beauty, craftsmanship, or hidden loveliness. It changes battlefields from making to choosing, handwork to mental action. This important artistic shift accepts prankishness while insisting jokes can carry serious arguments.

But this is one party's testimony, not final judgment. Friendly advocacy naturally thins the opposition. Compare a century's explanations more neutrally next.

\section[White Ceramic, Three Explanations]{One Piece of White Ceramic, Three Explanations}
Separate \textbf{what happened}, submission, removal, resignation; \textbf{why}, motives and importance; \textbf{contemporary understanding}, organizers denying art and friends affirming new thought; and \textbf{our evaluation}, liberation or sophisticated scam. Testimony deserves recording, not automatic truth. We compare why.

\textbf{First: an institutional test.}

Duchamp designed an experiment: pseudonym, full payment, correct procedures, testing the no-jury promise. The object was an examination institutions failed. Resignation rather than argument fits an experimenter recording results; anonymity controls celebrity bias. Yet this reduces Fountain to a one-off news event. Once tested, why does the prop still hold viewers a century later?

\textbf{Second: an aesthetic relocation.}

Selection extracts utility and forces attention to ceramic curves, symmetry, sheen. A Blind Man writer compared its reclining shape to a seated Buddha. Duchamp later described choosing readymades through visual indifference, neither beauty nor ugliness, forcing reasons beyond beauty. Yet interpreting a barely suppressed prank as formal appreciation feels like retrospective ennoblement, a solemn résumé appended to a joke.

\textbf{Third: a symptom of its age.}

In 1917, European powers had killed under civilization's banner for nearly three years, using civilization's latest machine guns and gas. Zurich Dada artists mocked high culture with nonsense poetry, collage, and noise. If civilization produced trenches, its aesthetics deserved no worship. Fountain can be this demolition's best-known blow. Broadest in scope, this is riskiest: Duchamp resisted heavy meanings, discussing the work coolly and evasively. Making him the age's accuser may hand him our own script.

Each grasps part and misses part. This is not indecision but method: an explanation covering everything may have discarded something beforehand.

\section{A Deeper Challenge: The Invisible Difference}
Introduce our weightiest theoretical device by advancing to 1964.

At a New York gallery, Warhol exhibited plywood Brillo Boxes silk-screened with commercial designs, nearly indistinguishable beside actual detergent cartons from warehouses. Arthur Danto visited and wrote "The Artworld." If eyes cannot distinguish them, the difference lies beyond sight: theoretical atmosphere, art-historical knowledge, people and institutions treating the object as a work. This invisible world makes the gallery box art; the supermarket carton lacks its discourse.

George Dickie refined institutional theory: artworld practices grant an object candidate status for appreciation. "Institution" simply means shared rules. A ball crossing a white line in a field means little; on a World Cup pitch, rules, referees, and spectators make it a goal. A museum relates to a urinal as the pitch to the ball.

Powerful explanation, substantial controversy: three accounts remain outstanding.

First, circularity. Art defines the artworld, which defines art. Is chicken-and-egg an explanation or a tongue-twister?

Second, gatekeepers. Curators, critics, dealers, collectors, a small group holding taste and capital, decide. The theory locates power precisely but does not justify its use. It explains entry, not whether entry ought to occur.

Third, outsiders. Children's drawings, remote craftspeople unfamiliar with artworlds, beautiful utilitarian objects: must all lose artistic citizenship? The theory risks excluding most human beauty-making.

Now read the crucial distinction carefully: \textbf{anything becoming art does not mean every judgment lacks reasons}.

Even if correct, institutional theory answers classification, not evaluation. Gallery admission assures art status, not goodness. Novelty, historical depth, skillful handling, honesty remain arguable reasons. Remove the entrance threshold and the evaluative court still sits.

A counterexample fixes this. Dutch painter Han van Meegeren forged Vermeers in the early twentieth century, fooling leading experts, one work hailed among Vermeer's greatest. After the war, accused of selling a national treasure to Nazi leader Göring, he defended himself: it was his forgery. Under supervision he painted another "Vermeer" to prove it. Once exposed, acclaimed masterpieces became forgeries in another category. Canvas and pigment remained; provenance, context, changed.

Some misuse this as proof experts pretend and judgment is arbitrary. Instead, it shows judgment \textbf{has} reasons: authorship, date, historical position. Experts were deceived by false reasons, not acting without reasons. Forgery's success presupposes reasons; if judgment were dice-throwing, why forge?

The other half cuts equally: carrying your toilet into a museum does not make you another Duchamp. A first introduction of the ordinary can occur only once in 1917. Time belongs to context. Repetition quotes rather than originates challenge. Readymades' pioneering value partly rests on its unrepeatability.

\section{Today: Museum Walls in Your Pocket}
The century-old question occurs daily with new props.

A banana first. At a major Miami art fair in 2019, Maurizio Cattelan taped a real banana to a wall and titled it Comedian. Editions reportedly sold for about 120,000 dollars each, really certificates and rights to install a banana, replaceable when rotten. Another performance artist removed and ate it as "Hungry Artist." A replacement continued the display. Every 1917 component remains: readymade, naming, institution, public, provocation, dispute, faster and dizzyingly dearer. You may see clever quotation or the fate Duchamp feared: institutional challenge becoming the institution's costliest commodity.

Your phone next. The same screenshot in different chats with different captions reverses meaning. You practice context changing meaning daily without certificates. Memes are popular readymades: borrowed images, added words, transformed significance. Every meme-maker knows selection, appropriation, renaming. In 1917 it required challenging the artworld; now a thumb suffices.

Museums move oppositely, collecting ration coupons, enamel mugs, old phones, particular years' masks and entry passes. Users never thought them exhibits. Display cases, context, turn them into evidence. The question becomes when ordinary objects silently become exhibits, often when they disappear.

AI is the newest battlefield. A prompt produces an image: who authors it? If selection outweighs manufacture, as Duchamp held, how does prompting, choosing one of a hundred results, signing, and publishing differ from R. Mutt? If not different, reconsider Fountain's discomfort; if different, identify the difference. Your generation needs it. The argument has barely begun, with the 1917 urinal watching from the witness stand.

\section{Your Question: An Empty Snack Bag}
Return to your school.

The art festival promises no jury, every registration displayed. The 1917 pledge returns to your corridor.

A classmate who sleeps through lessons submits an empty crisp packet, unwashed, stapled shut, titled Third Period This Afternoon, demanding the center of the exhibition.

You are on the committee. Display it or not?

Weigh the evidence before deciding.

For display: promises matter; the Independents broke theirs precisely here, a secret jury undoing the exhibition's meaning. Nonmanufacture does not mean no content; Duchamp, scholar's rocks, and tea bowls support him. Refusal as non-art repeats the gesture of jurors rejecting unfamiliar painting whom you dislike.

Against: broad classification does not erase evaluation. It can be art and lazy art. Duchamp's firstness mattered; copied provocation may provoke without challenging. Shared walls and everyone's time still need care, a duty not abolished by no-jury rules but passed to every organizer.

A deeper question remains yours: if anything can become art, what use remains for the word? Does a category containing everything explain nothing? Conversely, close the door tightly and how will unforeseen forms enter? Film, photography, and graffiti all once waited outside for a committee's admission.

There is no standard answer. But defending or rejecting that bag makes you more than a gallery spectator. Beside the white urinal, you join the newest committee in a century-long argument.
